\documentclass[a4paper,12pt]{article}
\usepackage[utf8]{inputenc}
\usepackage[italian]{babel}
\usepackage{listings}
\usepackage{xcolor}
\usepackage{float}
\usepackage{placeins}

\lstset{
  language=C,
  backgroundcolor=\color{lightgray!20},
  frame=single,
  breaklines=true,
  numbers=left,
  numberstyle=\tiny,
  stepnumber=1,
  numbersep=5pt,
  captionpos=b,
  basicstyle=\ttfamily\small,
  keywordstyle=\color{blue}\bfseries,
  commentstyle=\color{green!60!black},
  stringstyle=\color{red},
  showstringspaces=false,
  float=H
}

\title{LabReSiD26 -- Hands-On 5}
\author{Davide Ferrara -- 518629}
\date{}

\begin{document}

\maketitle

% ---------------------------------------------------------------------------
\section{Esercizio 1: Analisi di \texttt{server\_select.c}}
% ---------------------------------------------------------------------------

\subsection*{Domanda 1 -- Differenza tra \texttt{server\_fd} e \texttt{client\_fd}}

\textbf{server\_fd} è il socket di ascolto,
il suo unico scopo è ricevere le richieste di connesione.
Quando \texttt{select()} segnala che \texttt{server\_fd} è pronto,
significa che un nuovo client vuole connettersi --- si chiama \texttt{accept()}.

Il socket di connessione (\texttt{client\_fd}) viene creato da \texttt{accept()} ogni volta
che un client si connette. È il canale dedicato a quella specifica connessione: su di esso
si usano \texttt{recv()} e \texttt{send()}. Ne esiste uno per ogni client connesso
contemporaneamente.

\subsection*{Domanda 2 -- Perché \texttt{client\_fds[]} inizializzato a \texttt{-1} e non \texttt{0}}

I file descriptor partono da 0: \texttt{fd=0} è \texttt{stdin}, un fd valido.
Inizializzare \texttt{client\_fds[]} a \texttt{0} avrebbe due conseguenze errate:
\begin{itemize}
  \item Ogni slot risulterebbe già occupato (\texttt{client\_fds[i] != -1} sarebbe vero),
        impedendo l'inserimento di qualsiasi client reale, come dimostrato 
        dal seguente output:
        \begin{lstlisting}
>>> troppi client, rifiuto fd=4
        \end{lstlisting}
  \item \texttt{FD\_SET(0, \&read\_set)} aggiungerebbe \texttt{stdin} al set monitorato da
        \texttt{select()}, producendo notifiche spurie ogni volta che l'utente scrive da
        tastiera.
\end{itemize}
Con \texttt{-1} nessun fd reale assume mai quel valore, quindi ogni slot con
\texttt{client\_fds[i] == -1} è libero.

\subsection*{Domanda 3 -- Perché \texttt{max\_fd} parte da \texttt{server\_fd} e non da \texttt{0}}

\texttt{select(nfds, ...)} monitora i file descriptor da 0 fino a \texttt{nfds-1}.
Il primo argomento deve quindi essere il fd più grande tra quelli monitorati, più uno.
All'avvio l'unico fd presente è \texttt{server\_fd} (nel nostro caso \texttt{fd=3}),
quindi \texttt{max\_fd = server\_fd} e la chiamata diventa \texttt{select(4, ...)}.

Se si inizializzasse \texttt{max\_fd = 0}, la chiamata diventerebbe \texttt{select(1, ...)}:
il kernel scandisce il range \texttt{[0, 0]}, ovvero solo \texttt{fd=0} (stdin), e ignora completamente
\texttt{server\_fd=3}. Il server si bloccherebbe su \texttt{select()} indefinitamente
anche se un client tenta di connettersi come verificato sperimentalmente, in cui
il server non produce alcun output nonostante le connessioni in arrivo.

\subsection*{Domanda 4 -- Perché \texttt{FD\_ZERO} e \texttt{FD\_SET} ad ogni iterazione}

\texttt{select()} sovrascrive il set in input lasciando solo i fd pronti in quel ciclo.
Per non perdere gli fd non serviti, il set va ricostruito ad ogni iterazione usando
\texttt{client\_fds[]} come rubrica di verità.

\subsection*{Domanda 5 -- Disconnessione del secondo client}

Il server non crasha. Quando \texttt{recv()} ritorna 0 (connessione chiusa dal client),
il server chiude il fd con \texttt{close()}, imposta lo slot a \texttt{-1} liberandolo,
e ricalcola \texttt{max\_fd} scorrendo tutta la rubrica.
Gli altri client continuano a funzionare normalmente.

Output osservato (tre client connessi, terzo disconnesso con CTRL+C):
\begin{lstlisting}
[+] nuovo client connesso  fd=4
[+] nuovo client connesso  fd=5
[+] nuovo client connesso  fd=6
[-] client disconnesso (fd=6)
[fd=4] messaggio: "ciao" (4 byte)
[fd=5] messaggio: "ciao" (4 byte)
\end{lstlisting}

\subsection*{Domanda 6 -- \texttt{select()} può ritornare con \texttt{n > 1}}

Sì, \texttt{select()} può ritornare con \texttt{n > 1} quando più fd diventano pronti
nello stesso istante. Lo script \texttt{test\_multifd.sh 5} lo dimostra: apre 5 connessioni,
attende 1 secondo, poi tutti e 5 i client inviano il messaggio contemporaneamente.
L'output mostra prima 5 cicli con \texttt{n=1} (una connessione alla volta) e poi un ciclo
con \texttt{n=5} (tutti i messaggi arrivano insieme):

\begin{lstlisting}
>>> select() ritorna: 1 fd pronti   // client 1 si connette
>>> select() ritorna: 1 fd pronti   // client 2 si connette
>>> select() ritorna: 1 fd pronti   // client 3 si connette
>>> select() ritorna: 1 fd pronti   // client 4 si connette
>>> select() ritorna: 1 fd pronti   // client 5 si connette
>>> select() ritorna: 5 fd pronti   // tutti mandano insieme
[fd=4] messaggio: "messaggio dal client 1" (22 byte)
[fd=5] messaggio: "messaggio dal client 2" (22 byte)
[fd=6] messaggio: "messaggio dal client 3" (22 byte)
[fd=7] messaggio: "messaggio dal client 4" (22 byte)
[fd=8] messaggio: "messaggio dal client 5" (22 byte)
\end{lstlisting}

Il valore di ritorno di \texttt{select()} dipende quindi da quanti fd sono pronti
al momento della chiamata, non è sempre 1.

Con \texttt{./test\_multifd.sh 20} il server smette di accettare connessioni oltre la
decima: i client dall'undicesimo in poi vengono rifiutati con il messaggio
\texttt{[!] troppi client}.
Questo non è un problema di \texttt{select()}, ma un limite della struttura dati:
\texttt{MAX\_CLIENTS = 10} fissa la capacità massima della rubrica \texttt{client\_fds[]}.
I fd oltre \texttt{fd=13} (i primi 10 slot occupano fd 4--13) vengono chiusi immediatamente
dopo \texttt{accept()}.

\end{document}
